\documentclass[12pt]{article}
\usepackage[utf8]{inputenc}
\usepackage[T1]{fontenc}
\usepackage[margin=1in]{geometry}
\usepackage{setspace}
\usepackage{xcolor}
\usepackage[colorlinks=true, urlcolor=blue!60!black, linkcolor=blue!60!black]{hyperref}
\setlength{\parindent}{0em}
\setlength{\parskip}{0.6em}
\renewcommand{\baselinestretch}{1.0}
\begin{document}
\begin{center}
{\LARGE \textit{United States v. Rahimi}, 602 U.S. 680 (2024)\par}
\vspace{0.5em}
{\large Case Brief --- Two-Page Edition\par}
\end{center}

\section{Memory Jogger}
Clarifying \textit{Bruen}, the Court upheld 18 U.S.C.\ \S\ 922(g)(8): a person found by a court to pose a credible threat to another's physical safety may be temporarily disarmed (\href{https://www.law.cornell.edu/supremecourt/text/22-915}{602 U.S. 680}).

\section{Detailed Case Facts}
Zackey Rahimi assaulted his girlfriend and fired a gun as she fled. A Texas court issued a protective order finding he posed ``a credible threat'' to her ``physical safety.'' He violated it, was linked to five shootings, and was found with a pistol and rifle at home.

\section{Procedural History}
Charged under \S\ 922(g)(8), Rahimi pleaded guilty after losing a facial challenge; after \textit{Bruen}, the Fifth Circuit vacated the conviction for lack of a historical analogue. The Supreme Court reversed.

\section{Judicial Votes}
Reversed, 8--1. Roberts wrote for the Court; Sotomayor (with Kagan), Gorsuch, Kavanaugh, Barrett, and Jackson concurred; Thomas dissented.

\section{Holding}
Section 922(g)(8)(C)(i) is constitutional as applied to Rahimi and survives his facial challenge. A modern law need not be a ``historical twin''; it must fit the principles of the regulatory tradition in ``why'' and ``how'' it burdens the right.

\section{Analysis of Opinions}
Roberts held that the precedents ``were not meant to suggest a law trapped in amber.'' Surety laws and ``going armed'' laws show that people who threaten violence may be disarmed, and \S\ 922(g)(8)(C)(i) imposes temporary disarmament after an individualized judicial finding. He rejected the Government's ``not responsible'' theory. Sotomayor warned against a ``one-way ratchet''; Gorsuch and Kavanaugh defended text and history over balancing; Barrett said ``[h]istorical regulations reveal a principle, not a mold''; Jackson saw ``little method to \textit{Bruen}'s madness.'' Thomas, dissenting, found that surety laws imposed a lesser burden and affray laws addressed public terror, so ``[n]ot a single historical regulation justifies'' the statute.

\section{Examples: Future Applications}
\textbf{Same-Side:} a red-flag order disarming a respondent after a hearing finding a credible threat to others is upheld. \textbf{Opposite-Side:} permanently disarming a nonviolent felon without any threat finding is doubtful. \textbf{Fence-Sitter:} a boilerplate order barring force without a threat finding, under the reserved \S\ 922(g)(8)(C)(ii).

\section{Critique}
Strict originalists call the decision ``\textit{Bruen}-lite'' and note the Justices' split over tradition. Progressive scholars such as Joseph Blocher and Reva Siegel say it left \textit{Bruen}'s level-of-generality problem unsolved and that Founding-era analogues privilege an era that excluded women.

\section{Key Quotations}
\begin{quote}
``Taken together, the surety and going armed laws confirm what common sense suggests: When an individual poses a clear threat of physical violence to another, the threatening individual may be disarmed.''
\end{quote}
\noindent --- Chief Justice Roberts, Opinion of the Court, slip op.\ at 13.

\end{document}
